\section{Every fixed order for the small deformation regime}
\label{spb:ep:sec:subcritical}
Throughout this section $p\ge2$ is a fixed integer. Thus
$\lambda\asymp_p n^{1/q}$ and
\begin{equation}\label{spb:ep:eq:small}
 \lambda^2/n\asymp_p n^{-(p-1)/q}\longrightarrow0.
\end{equation}
This is the small parameter that makes an ordinary exponential expansion
of the defect valid after Poisson averaging.

\subsection{Global domination and tails}
\begin{lemma}\label{spb:ep:lem:sign}
For every integer $0\le r\le n$, $\Delta_n(r)\le0$.
Consequently
\begin{equation}\label{spb:ep:eq:globalR}
 0<\Rcal_p(n)\le q\Pp(X\equiv n\pmod q)
             =1+O_p(e^{-\kappa_q\lambda}).
\end{equation}
\end{lemma}
\begin{proof}
Let
\[
 F(r)=\frac{pn+r}{q}\log(1+pr/n)-\frac{p^2r}{q}.
\]
Direct differentiation gives $F(0)=F'(0)=0$ and
\begin{equation}\label{spb:ep:eq:Fsecond}
 F''(r)=\frac{p(2-p^2+pr/n)}{qn(1+pr/n)^2}.
\end{equation}
For $0\le r\le n$, its numerator is at most
$p(2-p^2+p)\le0$. Hence $F(r)\le0$.
Every logarithm in the sum in \eqref{spb:ep:eq:Delta} is nonpositive.
This proves the sign. Now apply \eqref{spb:ep:eq:exactpoisson} and
\eqref{spb:ep:eq:filter} with $j=0$.
\end{proof}

The right side of \eqref{spb:ep:eq:globalR} need not be identically one:
residue probabilities may have exponentially small oscillations.
Nevertheless it already excludes a positive algebraic first correction.
In particular, a positive $n^{-1/4}$ coefficient for $p=3$ is impossible.

The Poisson moment generating function gives, with
$h=2\log2-1>0$,
\begin{equation}\label{spb:ep:eq:tail}
 \Pp(X>2\lambda)\le e^{-h\lambda}.
\end{equation}
Indeed, exponential Markov with parameter $\log2$ bounds the probability
by $\exp(\lambda(2-1)-2\lambda\log2)$.
For each fixed integer $d\ge0$, Cauchy--Schwarz yields
\begin{equation}\label{spb:ep:eq:weightedtail}
 \E[X^d\1_{\{X>2\lambda\}}]
   \le (\E X^{2d})^{1/2}\Pp(X>2\lambda)^{1/2}
   =O_d((1+\lambda)^d e^{-h\lambda/2}).
\end{equation}
By Lemma~\ref{spb:ep:lem:sign}, the exact normalized contribution outside
$X\le2\lambda$ is bounded by the same unweighted tail, up to $q$.
For large $n$, $2\lambda<n$, so the original cutoff at $n$ causes no
additional difficulty. These are global estimates, including the far
endpoint $r=n$.

\subsection{Defect polynomials and a uniform analytic remainder}
For $j\ge1$ define
\begin{equation}\label{spb:ep:eq:Dgeneral}
 D_j(r)=(-1)^{j+1}\left\{
  \frac{p^j}{q}\left(\frac1j-\frac{p^2}{j+1}\right)r^{j+1}
  +\frac1j\sum_{h=0}^{r-1}\bigl[(pr-qh)^j-(pr)^j\bigr]\right\}.
\end{equation}
The finite sum is interpreted as its polynomial in $r$, obtained for
example from Faulhaber sums or interpolation in the falling-factorial
basis. Thus $D_j\in\Q[r]$, $D_j(0)=0$, and
$\deg D_j\le j+1$.
Define
\begin{equation}\label{spb:ep:eq:Hrec}
 H_0(r)=1,\qquad
 H_j(r)=\frac1j\sum_{s=1}^j sD_s(r)H_{j-s}(r).
\end{equation}
Formal differentiation of an exponential shows that
\[
 \exp\left(\sum_{j\ge1}D_j(r)z^j\right)
        =\sum_{j\ge0}H_j(r)z^j.
\]
Induction gives $\deg H_j\le2j$ and $H_j(0)=0$ for $j\ge1$.

\begin{lemma}\label{spb:ep:lem:cauchy}
For each fixed $S\ge0$, uniformly over integers $0\le r\le2\lambda$,
\begin{equation}\label{spb:ep:eq:pointrem}
 e^{\Delta_n(r)}=\sum_{j=0}^S H_j(r)n^{-j}
       +O_{p,S}\bigl((r+1)^{2S+2}n^{-S-1}\bigr).
\end{equation}
\end{lemma}
\begin{proof}
For an integer $r\ge1$ consider the analytic germ
\begin{equation}\label{spb:ep:eq:G}
 G_r(z)=e^{-p^2r/q}
       (1+prz)^{p/(qz)-pr/q}
       \prod_{h=0}^{r-1}(1+(pr-qh)z),
\end{equation}
using logarithmic branches at $z=0$. The apparent singularity at zero is
removable and $G_r(0)=1$. Recombining the denominators in
\eqref{spb:ep:eq:Delta} gives $G_r(1/n)=e^{\Delta_n(r)}$.

Choose a small constant $\eta_p>0$ and put
$\rho_r=\eta_p/(r+1)^2$. For $|z|\le\rho_r$, every logarithm argument in
\eqref{spb:ep:eq:G} lies in a fixed disk about one excluding zero. After its
constant $p^2r/q$ is subtracted, the expression
$(p/(qz))\log(1+prz)$ is $O_p(r^2|z|)$.
The remaining logarithm term and the sum of the $r$ product logarithms
are also $O_p(r^2|z|)$. Therefore
\[
 |\log G_r(z)|\le C_p r^2|z|\le C'_p,
 \qquad |G_r(z)|\le e^{C'_p},
\]
uniformly in $r$.
Expanding its logarithm gives exactly \eqref{spb:ep:eq:Dgeneral}; exponentiation
gives \eqref{spb:ep:eq:Hrec}.

Cauchy's coefficient estimate on $|z|=\rho_r$ and a geometric tail give
\[
 \left|G_r(1/n)-\sum_{j=0}^S H_j(r)n^{-j}\right|
  \le C_{p,S}(r+1)^{2S+2}n^{-S-1}
\]
whenever $(r+1)^2/n\le\eta_p/2$.
This condition holds uniformly for $r\le2\lambda$ by \eqref{spb:ep:eq:small}.
For $r=0$, $G_0=1$ and the conclusion is immediate.
\end{proof}

\subsection{The coefficient theorem}
Define the Poisson transform of $H_j$ by
\begin{equation}\label{spb:ep:eq:U}
 H_j(r)=\sum_d h_{j,d}r^d,
 \qquad U_j(z)=\sum_d h_{j,d}T_d(z).
\end{equation}
Thus $U_0=1$, $U_j(0)=0$ for $j\ge1$, and $\deg U_j\le2j$.

\begin{theorem}\label{spb:ep:thm:allp}
For every fixed $S\ge0$,
\begin{equation}\label{spb:ep:eq:allH}
 \Rcal_p(n)=\sum_{j=0}^S n^{-j}U_j(\lambda)
   +O_{p,S}\left(\frac{(1+\lambda)^{2S+2}}{n^{S+1}}
                              +e^{-c_{p,S}\lambda}\right)
\end{equation}
for some $c_{p,S}>0$. In particular the exponential term can be absorbed
into the algebraic remainder. Setting $t=n^{-1/q}$ and $\lambda=c_p/t$,
for every fixed $K\ge0$ this implies \eqref{spb:ep:eq:overviewfixed}, with
\begin{equation}\label{spb:ep:eq:Cextract}
 C_{p,k}(c)=[t^k]\sum_{j=0}^{\lfloor K/(p-1)\rfloor}
                         t^{qj}U_j(c/t),\qquad 0\le k\le K.
\end{equation}
\end{theorem}
\begin{proof}
Apply \eqref{spb:ep:eq:pointrem} in the exact expectation on $X\le2\lambda$.
The expectation of the remainder is at most a constant times
$n^{-S-1}\E(X+1)^{2S+2}$, giving the algebraic term in \eqref{spb:ep:eq:allH}.
The exact tail is controlled by \eqref{spb:ep:eq:tail} and the sign lemma.
The tails of the finitely many polynomials $H_j$ are controlled by
\eqref{spb:ep:eq:weightedtail}. Remove their residue restrictions using
\eqref{spb:ep:eq:filter}. All resulting exponential terms, including fixed
polynomial factors, can be combined into $O(e^{-c_{p,S}\lambda})$.

A monomial $n^{-j}\lambda^d$ becomes $c^dt^{qj-d}$, with
$1\le d\le2j$ when $j\ge1$. Its possible exponent lies in
\begin{equation}\label{spb:ep:eq:bands}
 [(p-1)j,\ qj-1].
\end{equation}
Take $S=\lfloor K/(p-1)\rfloor$.
The remainder in \eqref{spb:ep:eq:allH} starts at order
$t^{(p-1)(S+1)}=O(t^{K+1})$.
Discarding finitely many polynomial terms with exponent larger than $K$
also gives $O(t^{K+1})$. This proves \eqref{spb:ep:eq:Cextract} and the theorem.
\end{proof}

Formula \eqref{spb:ep:eq:bands} gives a structural exclusion of low powers.
It does not imply that every exponent inside an allowed band has a
nonzero coefficient. In particular, the absence of
$t,t^2,\ldots,t^{p-2}$ requires no symbolic computation.

\subsection{Explicit corrected terms}
The first two defect polynomials are
\begin{align}
 D_1(r)&=-Ar^2+\frac q2r,\label{spb:ep:eq:D1}\\
 D_2(r)&=\gamma_3r^3-\gamma_2r^2-\gamma_1r,\label{spb:ep:eq:D2}\\
 \gamma_3&=\frac{p^3}{3}-\frac1{6q},\qquad
 \gamma_2=\frac{p^2-1}{4},\qquad \gamma_1=\frac{q^2}{12}.\nonumber
\end{align}
Since $U_1(\lambda)=\E D_1(X)=-A\lambda^2-b\lambda$,
Theorem~\ref{spb:ep:thm:allp} gives \eqref{spb:ep:eq:firstuniv}.
For $p=2$, write $c=e^{4/3}/3^{1/3}$ and $t=n^{-1/3}$. Then
\begin{align}\label{spb:ep:eq:p2}
 \Rcal_2(n)={}&1-\frac{13c^2}{6}t
 +\left(\frac{169c^4}{72}-\frac{2c}{3}\right)t^2\nonumber\\
 &+\left(\frac{121c^3}{9}-\frac{2197c^6}{1296}\right)t^3
 +O(t^4).
\end{align}
One additional coefficient is
\[
 C_{2,4}(c)=\frac{134c^2}{9}-\frac{2977c^5}{108}
                         +\frac{28561c^8}{31104}.
\]
For $p=3$, write $c=e^{9/4}/\sqrt2$ and $t=n^{-1/4}$. Then
\begin{align}\label{spb:ep:eq:p3}
 \Rcal_3(n)={}&1-\frac{37c^2}{8}t^2-\frac{21c}{8}t^3
        +\frac{1369c^4}{128}t^4+O(t^5)\nonumber\\
 ={}&1-\frac{37e^{9/2}}{16}n^{-1/2}
       -\frac{21e^{9/4}}{8\sqrt2}n^{-3/4}
       +\frac{1369e^9}{512}n^{-1}+O(n^{-5/4}).
\end{align}
The next coefficients are
\[
 C_{3,5}(c)=\frac{12265c^3}{192},\qquad
 C_{3,6}(c)=\frac{9471c^2}{128}-\frac{50653c^6}{3072}.
\]
All these identities are rational-polynomial outputs of the finite
recurrence, not numerical coefficient fits.

\begin{remark}[Formal order and useful numerical accuracy]
The theorem is asymptotic for each fixed truncation. It does not say that
larger order improves every finite-$n$ approximation. The constants in
\eqref{spb:ep:eq:p3} are large. At $n=10000$, even the fractional-power
truncation through $t^{12}$ has relative error about $-0.662$.
Section~\ref{spb:ep:sec:numerics} gives the complete-sum diagnostic and contrasts
it with another legitimate finite truncation of \eqref{spb:ep:eq:allH}.
\end{remark}
